\documentclass[11pt]{article}

\usepackage[margin=1in]{geometry}
\usepackage{titlesec}
\usepackage{enumitem}
\usepackage{listings}
\usepackage{xcolor}
\usepackage{hyperref}

\hypersetup{colorlinks=true, linkcolor=blue, urlcolor=blue}

\definecolor{codebg}{rgb}{0.96,0.96,0.96}
\lstset{
  basicstyle=\ttfamily\small,
  backgroundcolor=\color{codebg},
  breaklines=true,
  frame=single,
  columns=fullflexible,
  keepspaces=true,
}

\titleformat{\section}{\large\bfseries}{\thesection}{1em}{}
\titleformat{\subsection}{\normalsize\bfseries}{\thesubsection}{1em}{}

\title{Density-Based Traffic Signal Control \\ \large Change Log: Broker Migration, Emergency Preemption, and RTX~2060 Performance Tuning}
\author{Engineering note}
\date{2026-06-15}

\begin{document}
\maketitle

\section{Summary}
This document records the changes made to the density-based traffic signal
control system in the \texttt{tl} repository. The work covers four areas:

\begin{enumerate}[nosep]
  \item Migrating the MQTT broker from \texttt{test.mosquitto.org} to
        \texttt{broker.hivemq.com} across the firmware and documentation.
  \item Confirming and documenting the two-lane, density-driven signal logic
        (one light for the vertical lane, one for the horizontal lane).
  \item Adding \emph{emergency-vehicle preemption}: an ambulance or firetruck
        forces its lane green, but only once we are confident the detection is
        real.
  \item Performance optimizations targeting a client GPU no larger than an
        NVIDIA RTX~2060 running the smallest YOLO model, while keeping both the
        live-camera (RTSP) and video-file pipelines intact.
\end{enumerate}

\section{System Context}
The pipeline captures video from four IP cameras (or local video files for
testing), runs a custom-trained YOLO11 detector
(classes \texttt{ambulance}, \texttt{firetruck}, \texttt{vehicle}), counts
vehicles per direction, and decides which of the two traffic lights should be
green. Decisions are published over MQTT to topics
\texttt{ub-traffic-light/signals/vertical} and
\texttt{.../horizontal}. An ESP32 subscribed to each topic drives the physical
lights and coordinates the partner light over ESP-NOW.

\section{1. MQTT Broker Migration to HiveMQ}
The Python publisher (\texttt{gui/publisher.py}) and the manual test scripts
already pointed at \texttt{broker.hivemq.com}. The remaining references to
\texttt{test.mosquitto.org} were updated so the whole system is consistent:

\begin{itemize}[nosep]
  \item \texttt{esp/bidirectional\_mqtt/bidirectional\_mqtt.ino}
  \item \texttt{esp/bidirectional\_mqtt\_vertical/bidirectional\_mqtt\_vertical.ino}
  \item \texttt{esp/bidirectional\_mqtt\_horizontal/bidirectional\_mqtt\_horizontal.ino}
  \item \texttt{install.txt} (the \texttt{MQTT\_SERVER} configuration line)
\end{itemize}

\noindent Each \texttt{const char* MQTT\_SERVER} now reads
\texttt{"broker.hivemq.com"} on port \texttt{1883}. No topic names or client-ID
logic changed, so the existing firmware behaviour is preserved.

\section{2. Two-Lane Density Logic}
The intersection has exactly two phases. Cameras are tagged with a
\texttt{direction} of \texttt{vertical} or \texttt{horizontal}; collinear lanes
share a phase. Per inference cycle the system sums the \texttt{vehicle}
detections for each direction and prefers to give green to the busier
direction:

\begin{lstlisting}
if vertical_count > horizontal_count:  -> green: vertical
elif horizontal_count > vertical_count: -> green: horizontal
else:                                    -> hold current phase
\end{lstlisting}

\noindent Two timing guards keep the lights sane and safe:
\begin{itemize}[nosep]
  \item \textbf{Minimum green} (\texttt{MIN\_GREEN}, 15\,s): the system will not
        switch away from a freshly granted green until this elapses, preventing
        flicker when counts are close.
  \item \textbf{Maximum green} (\texttt{MAX\_GREEN}, 85\,s): if one direction has
        held green this long, the opposite direction is forced green to avoid
        starvation.
\end{itemize}

\noindent Every phase change is mediated by a 3\,s yellow transition on the
lane losing priority, after which its partner light is told to go green.

\section{3. Emergency-Vehicle Preemption}
\subsection{Goal}
If an ambulance or firetruck is present in a lane, that lane should turn green
immediately. However, a single noisy frame must \emph{not} hijack the
intersection, so a detection only counts as an emergency once we are confident
it is genuine.

\subsection{Confidence gate}
A detection is treated as a candidate emergency only if its class is
\texttt{ambulance} or \texttt{firetruck} \emph{and} its confidence is at least
\texttt{EMERGENCY\_CONFIDENCE} (default \texttt{0.55}, higher than the
\texttt{0.30} used for ordinary vehicle counting).

\subsection{Persistence gate}
Per cycle, each lane's count of confident emergency detections updates a
\emph{streak} counter: the streak increments while emergencies are seen and
resets to zero the moment a lane shows none. A lane is only declared a
confirmed emergency once its streak reaches \texttt{EMERGENCY\_FRAMES}
(default \texttt{5} consecutive cycles). This rejects isolated false positives
while still reacting within a fraction of a second at typical frame rates.

\subsection{Preemption behaviour}
When a lane is confirmed:
\begin{itemize}[nosep]
  \item It is switched green immediately, \textbf{bypassing the minimum-green
        hold} on the current phase.
  \item While its emergency remains active it is \textbf{exempt from the
        maximum-green timeout}, so it stays green until the emergency vehicle
        clears.
  \item The 3\,s yellow transition on the losing lane is still honoured for
        safety.
  \item If both lanes show confirmed emergencies simultaneously, the currently
        green lane is held; otherwise the lane with more emergency detections
        wins.
\end{itemize}

\subsection{Tuning}
Both thresholds are configurable per deployment via the
\texttt{experimental\_settings} block in the camera config files
(\texttt{emergency\_confidence}, \texttt{emergency\_frames}) and are applied
live when settings change. The detector classes \texttt{ambulance} and
\texttt{firetruck} can be enabled in the on-screen bounding-box overlay for
visual verification.

\section{4. Performance Optimizations (RTX~2060, Smallest Model)}
The clients run on a GPU no larger than an RTX~2060 with the smallest YOLO
weights, so the goal was to cut per-frame overhead without changing detection
behaviour.

\subsection{Eliminated redundant GPU$\rightarrow$CPU transfers}
Previously each frame's detections were pulled off the GPU as separate
\texttt{xyxy}, \texttt{conf}, and \texttt{cls} tensors in the display path, the
counting path performed its own \texttt{cls} transfer, and the inference helper
computed shifted coordinates that were then \emph{discarded}. The detections
are now transferred \textbf{once} per frame as a single
\texttt{(N,6)} array \texttt{[x1,y1,x2,y2,conf,cls]}, already shifted into
full-frame coordinates, and cached. Both the signal logic and the renderer read
this cache. The dead transfer was removed entirely. On four cameras this
removes several host/device synchronisations per cycle, the most costly
non-inference work in the loop.

\subsection{cuDNN autotuning, TF32, and FP16}
Because the inference tile size is fixed frame-to-frame, cuDNN benchmarking is
enabled so it caches the fastest convolution kernels. TF32 matmul/convolution
paths are enabled, and the model is fused and run in FP16 on the GPU. These are
low-risk wins well suited to Turing-class hardware.

\subsection{Warm-up inference}
A single dummy inference is run at startup so the first real frame no longer
pays the one-off kernel-compilation and allocator stall. This also makes the
reported latency reflect steady-state performance.

\subsection{Correct per-camera ROI indexing}
The region-of-interest cropping now uses each frame's true camera index rather
than its position within the (possibly shorter) batch. This fixes a latent bug
where a dropped camera frame would apply the wrong ROI, and keeps inference
confined to the relevant road area.

\subsection{Shared pipeline path (\emph{the other implementation, kept})}
The repository has two capture paths: live RTSP cameras and local video files
(used for testing/demo). Both were retained. Their per-cycle processing
(batching, inference, signal update, rendering) was de-duplicated into a single
\texttt{\_process\_frames} helper, so both paths now share the identical,
optimized code and cannot drift apart. The video-file path additionally
benefits from all the optimizations above.

\section{Files Changed}
\begin{itemize}[nosep]
  \item \texttt{gui/inference.py} --- emergency preemption, single-transfer
        detection cache, cuDNN/TF32/FP16 setup, warm-up, ROI index fix, shared
        \texttt{\_process\_frames} helper, live-tunable thresholds.
  \item \texttt{gui/publisher.py} --- already on HiveMQ (no change required).
  \item \texttt{esp/*/*.ino} --- broker set to \texttt{broker.hivemq.com}.
  \item \texttt{configs/*.json} --- added \texttt{emergency\_confidence} and
        \texttt{emergency\_frames} to \texttt{experimental\_settings}.
  \item \texttt{install.txt} --- broker documentation updated.
\end{itemize}

\section{Validation Notes}
Python modules were syntax-checked and the configuration files were validated
as JSON. End-to-end behaviour (live detection latency on the RTX~2060 and the
emergency override against real ambulance/firetruck footage) should be
confirmed on the target hardware, and \texttt{EMERGENCY\_CONFIDENCE} /
\texttt{EMERGENCY\_FRAMES} tuned to the deployed camera frame rate.

\end{document}
